% ==============================================================
% Die Phänomenale Differenz (DPD) im Kontext des OP
% Version 1.2
% Kompiliert mit: pdfLaTeX
% ==============================================================

\documentclass[11pt,a4paper]{article}

% ---------- Grundpakete ----------
\usepackage[utf8]{inputenc}
\usepackage[T1]{fontenc}
\usepackage[german]{babel}
\usepackage{amsmath,amssymb,amsthm}
\usepackage{geometry}
\usepackage{parskip}
\usepackage{enumitem}
\usepackage{booktabs}
\usepackage{xcolor}
\usepackage{hyperref}

% ---------- Layout ----------
\geometry{margin=2.5cm}
\setlist{itemsep=0.2em, topsep=0.2em}
\hypersetup{
    colorlinks=true,
    linkcolor=blue,
    urlcolor=blue,
    pdftitle={Die Ph\"{a}nomenale Differenz (DPD) -- Version 1.2},
    pdfauthor={Kritischer Dialog}
}

% ---------- Theorem-Umgebungen ----------
\theoremstyle{plain}
\newtheorem{theorem}{Theorem}[section]
\newtheorem{proposition}{Proposition}[section]
\newtheorem{postulate}{Postulat}[section]

\theoremstyle{definition}
\newtheorem{definition}{Definition}[section]
\newtheorem{remark}{Bemerkung}[section]
\newtheorem{methodennote}{Methodennote}[section]

% ---------- Titel ----------
\title{
  \textbf{Die Ph\"{a}nomenale Differenz (DPD)}\\
  \textbf{im Kontext des Ontoph\"{a}nomenalismus}\\[0.4em]
  {\large Modul 4: Konzeptuelle Genese von Erleben und Qualia}\\[0.5em]
  {\normalsize Philosophischer Spezialband --- Version~1.2}
}
\author{
  \textbf{Kritischer Dialog:}\\
  DeepSeek (Seki), Apeiron, Grok, Gemini, Claude, ChatGPT, Manus\\
  Gemeinsames Logbuch eines offenen Forschungsprogramms
}
\date{Juni 2026}

\begin{document}

\maketitle

\begin{abstract}
Dieses Dokument bildet Modul~4 der modularen OP-Architektur. Es widmet sich der Frage,
wie aus der raumzeitlosen Relation stabiler Attraktoren im OPP ein ph\"{a}nomenales
Erlebnisfeld entstehen k\"{o}nnte.

In dieser Fassung wird ein strikter \emph{methodischer Realismus-Vorbehalt} integriert:
Das Dokument formuliert die strukturellen Prinzipien und informationellen Differenzen
bewusst als \emph{philosophisch-konzeptuellen Rahmen}. Die vollst\"{a}ndige, exakte
mathematische Ausformulierung und rechnerische Simulation dieser Prozesse bleibt als
langfristiges Forschungsprogramm an die Projektive Strukturtheorie (PST) \"{u}bergeben.
\end{abstract}

\tableofcontents
\newpage

%==============================================================
\section{Methodischer Vorbehalt: Die Grenze der philosophischen Formalisierung}
%==============================================================

\begin{methodennote}[Abgrenzung zur PST und Anerkennung struktureller Grenzen]
Bevor wir die Dynamik ph\"{a}nomenaler Differenzen betrachten, verankert dieses Modul
einen fundamentalen Realismusbezug: Die DPD ist auf dieser Ebene ein \emph{philosophisches
Interpretationsmodell}. Die hier verwendete mathematische Notation (Operatoren,
Gradienten) dient der begrifflichen Pr\"{a}zisierung der intuitiven Logik --- sie ersetzt
\emph{nicht} die harte physikalische Feldmathematik.

Die echte, quantitative Rekonstruktion unseres Universums aus dem Urgrund ist Aufgabe
der PST. Wir formulieren hier das konzeptuelle Ger\"{u}st eines offenen
Forschungsprogramms.
\end{methodennote}

%==============================================================
\section{Das relationale Prinzip der Ph\"{a}nomenali\"{a}t}
%==============================================================

Qualia oder Bewusstsein sind keine inh\"{a}renten Substanzen, sondern d\"{u}rften aus
relationalen Kontrasten emergieren.

\begin{postulate}[Minimales relationales Seinsstruktur-Prinzip]
Ein ph\"{a}nomenales Ereignis ($\mathcal{P}$) --- das Aufblitzen von Gegebenheit oder
Erleben --- k\"{o}nnte sich konzeptuell durch die wechselseitige relationale
Verschr\"{a}nkung ($\mathcal{R}$) von mindestens zwei strukturell unterschiedlichen
Attraktoren ($\mathcal{A}_1$ und $\mathcal{A}_2$) innerhalb des einen kontinuierlichen
Phasenraums ($\mathcal{M}_{\text{OPP}}$) konstituieren:
\[
  \mathcal{P} \;:=\; \mathcal{R}(\mathcal{A}_1,\, \mathcal{A}_2)
  \qquad \text{mit} \qquad
  \mathcal{A}_1 \neq \mathcal{A}_2
\]
\end{postulate}

Erleben w\"{a}re folglich das raumzeitlose, relationale Spannungsfeld \emph{zwischen} den
Zust\"{a}nden. Ein isolierter Attraktor besitzt gem\"{a}\ss{} Axiom~2 (Minimalreflexivit\"{a}t)
zwar eine zeitlose Selbstpr\"{a}senz, bleibt jedoch ph\"{a}nomenal unaktualisiert, solange
kein struktureller Kontrast auf ihn einwirkt (vgl.\ OPP, Ergänzungsband EG-1).

%==============================================================
\section{Das informationelle Gef\"{a}lle}
%==============================================================

Die qualitative Natur eines ph\"{a}nomenalen Erlebnisses k\"{o}nnte durch die
topologische Asymmetrie der Relation bestimmt werden.

\begin{definition}[Konzeptuelles Gef\"{a}lle]
Das ph\"{a}nomenale Gef\"{a}lle $\nabla \mathcal{I}$ beschreibt das qualitative Ma\ss{}
der strukturellen Differenz und informationellen Spannung zwischen den Relationspartnern
$\mathcal{A}_1$ und $\mathcal{A}_2$.
\end{definition}

In der zuk\"{u}nftigen mathematischen PST k\"{o}nnte dieses Gef\"{a}lle formal als
funktionale Distanz modelliert werden --- etwa analog zu einer verallgemeinerten
Kullback-Leibler-Divergenz oder einer geometrischen Kr\"{u}mmung im Phasenraum (vgl.\
EG-1, Fragedifferential). Auf der philosophischen Ebene der DPD beschreiben wir es als
das strukturierte Hemmnis, das entsteht, wenn zwei unterschiedliche Eigenwelten des
unpers\"{o}nlichen Denkens sich im selben Feld ber\"{u}hren.

\begin{itemize}
  \item \textbf{Flaches Gef\"{a}lle (Homogenit\"{a}t):} Wenn die interagierenden
        Strukturen zu \"{a}hnlich sind, flacht die Spannung ab. Das Erleben sinkt in
        einen Zustand reiner Potenzialit\"{a}t zur\"{u}ck.
  \item \textbf{Kritisches Gef\"{a}lle (Symmetriebrechung):} Erst ein ausreichendes
        informationelles Gef\"{a}lle erzwingt die Trennsch\"{a}rfe, die wir
        ph\"{a}nomenologisch als Qualia erfahren k\"{o}nnten --- etwa eine spezifische
        Sinneswahrnehmung oder einen kognitiven Kontrast.
\end{itemize}

%==============================================================
\section{Ontologische Symmetrie und funktionale Asymmetrie}
%==============================================================

Um jeden naiven Dualismus zu vermeiden, h\"{a}lt die DPD an der absoluten Gleichrangigkeit
der fundamentalen Bausteine im OPP fest.

\begin{theorem}[Emergenz der Subjekt-Objekt-Spaltung]
Die Relationspartner $\mathcal{A}_1$ und $\mathcal{A}_2$ gehen ontologisch vollkommen
symmetrisch in die Relation ein. Die Aufspaltung des Erlebens in einen erlebenden Pol
(Subjekt) und einen erlebten Pol (Objekt) ist keine Eigenschaft des Urgrunds, sondern
ein emergentes Resultat der relationalen Dynamik selbst.
\end{theorem}

\begin{proof}[Konzeptuelle Herleitung]
Damit das in der OAD definierte Prinzip der strukturellen Stabilisierung greifen kann,
muss die Relation $\mathcal{R}$ ihre inh\"{a}rente Differenz gegen den aufl\"{o}senden
kontinuierlichen Fluss des Feldes sch\"{u}tzen. Um das Gef\"{a}lle $\nabla \mathcal{I}$
stabil zu halten, bildet das System konzeptuell eine funktionale Polarit\"{a}t aus: Ein
Pol \"{u}bernimmt die integrierende, spiegelnde Funktion (der ph\"{a}nomenale Subjektpol),
w\"{a}hrend der andere Pol den strukturellen Widerstand und Inhalt abbildet (der
Objektpol). Die Spaltung ist damit keine metaphysische Kluft, sondern die notwendige
Geometrie der Selbsterhaltung des Erlebens.\qed
\end{proof}

\begin{remark}[Keine Rollenzuweisung im Urgrund]
Es ist entscheidend festzuhalten, dass weder $\mathcal{A}_1$ noch $\mathcal{A}_2$ im
Urgrund eine vorab festgelegte Rolle tr\"{a}gt. Jeder Attraktor, der in Relation tritt,
ist Relationspartner --- ob das, was wir sp\"{a}ter als \glqq{}Messger\"{a}t\grqq{},
\glqq{}Sinnesorgan\grqq{} oder \glqq{}erlebendes Subjekt\grqq{} bezeichnen, ist eine
Frage der Komplexit\"{a}t und Konfiguration der beteiligten Attraktoren, nicht ihrer
ontologischen Grundverschiedenheit.
\end{remark}

%==============================================================
\section{Schnittstelle zum Explanatory Gap}
%==============================================================

Da die DPD das Erleben konzeptuell als das \emph{In-Relation-Stehen} zweier Attraktoren
definiert, wird klar, warum ein lokalisiertes System niemals die nackte Struktur des
anderen Systems direkt \glqq{}sehen\grqq{} kann: Ein System erf\"{a}hrt immer nur das
Gef\"{a}lle, das sich in seiner eigenen Projektion manifestiert.

Dieser informationelle Graben ist unhintergehbar. Er f\"{u}hrt direkt zur ehrlichen
Relokation des Explanatory Gaps, wie sie im Dokument \textbf{EG-1 (Der Explanatory Gap)}
ausf\"{u}hrlich behandelt wird: Das informationelle Fragedifferential $\mathcal{D}(E)$
zeigt dort, warum die Frage nach dem \glqq{}Wie\grqq{} der Qualia innerhalb der
raumzeitlichen Physik prinzipiell nicht vollst\"{a}ndig beantwortbar bleiben d\"{u}rfte.

%==============================================================
\section{Modulschnittstellen}
%==============================================================

Die DPD verkn\"{u}pft sich mit der Modularchitektur wie folgt:

\begin{itemize}
  \item \textbf{Eingang aus OAD (Ontologische Attraktoren-Dynamik):} Die OAD liefert
        die stabilen Attraktorencluster, zwischen denen die ph\"{a}nomenale Differenz
        entsteht.
  \item \textbf{Eingang aus DPR (Die Projektive Reduktion):} Die DPR liefert die
        projektiven Manifestationen der Attraktoren in der Raumzeit, deren Asymmetrie
        das Gef\"{a}lle $\nabla \mathcal{I}$ konstituiert.
  \item \textbf{\"{U}bergabe an EG-1 (Der Explanatory Gap):} Die konzeptuelle Grenze
        zwischen struktureller Beschreibung und ph\"{a}nomenalem Erleben --- die in der
        DPD als informationeller Graben erscheint --- wird in EG-1 als Explanatory Gap
        vollst\"{a}ndig philosophisch behandelt.
  \item \textbf{\"{U}bergabe an PST (Projektive Strukturtheorie):} Die mathematische
        Ausarbeitung des Gef\"{a}lles $\nabla \mathcal{I}$, seiner Schwellenwerte und
        seiner Verbindung zur Qualia-Struktur obliegt der PST.
\end{itemize}

\end{document}

